\SourceLeaf{62}glückliche Erfolge bei drei weiteren, seither neu entdeckten Planeten vollständig ausgeräumt.

Jene Methoden bilden die ursprüngliche Grundlage des Werks, dem diese Anzeige gewidmet ist. Bei der fortgesetzten Beschäftigung mit dem Gegenstand hatte der Verfasser jedoch Gelegenheit, sie immer wieder so vielfach zu ergänzen, zu ändern und zu verbessern, dass von ihrer ursprünglichen Gestalt fast keine Spuren übrig geblieben sind. Er hofft daher, dass die Astronomen, die von Anfang an lebhaftes Interesse an diesen Untersuchungen und den Wunsch nach einer baldigen Veröffentlichung geäußert hatten, keinen Grund haben werden, mit dem späten Erscheinen unzufrieden zu sein. Wir überlassen den sachkundigen Lesern des Werks selbst das Urteil, inwieweit diese Hoffnung begründet ist, und beschränken uns hier darauf, seinen Aufbau und Inhalt in aller Kürze darzustellen.

Das mehrfach erwähnte Problem war zwar der Anlass und der Hauptzweck des vorliegenden Werks, konnte aber in der Ausführung nur einen Teil, ja nur den kleineren Teil davon ausmachen. Man kann sich leicht denken, dass sich eine angemessene Lösung nicht geben ließ, ohne eine Fülle bisher noch nicht entwickelter Wahrheiten über die Bewegung der Himmelskörper in Kegelschnitten vorauszusetzen. Diese neuen, schon an sich sehr interessanten Beziehungen waren so zahlreich und so eng miteinander verbunden, dass sie im Zusammenhang dargestellt werden mussten. Ihnen ist daher der gesamte erste Teil des Werks gewidmet. Der Aufbau ist sehr einfach. Dieser erste Teil gliedert sich in vier Abschnitte. Der erste und der zweite enthalten die Beziehungen, die sich auf einen einzigen Ort des Himmelskörpers auf seiner Bahn oder im Raum beziehen. Der dritte und der vierte entwickeln die Beziehungen zwischen zwei oder mehreren Orten. Neben den neuen Beziehungen mussten durchweg auch die wichtigsten bereits bekannten aufgeführt werden. Diese sind jedoch meist in einer dem Verfasser eigenen Form entwickelt. Der begrenzte Raum dieser Blätter erlaubt es nicht, auch nur einen Überblick über alle hier behandelten Gegenstände zu geben. Wir beschränken uns darauf, einige der umfangreicheren Untersuchungen zu berühren. Dazu gehört im ersten Abschnitt ein Exkurs
\SourceLeaf{63}\SourceGathering{VI.~8}über die Genauigkeit, die sich bei numerischen Rechnungen mit Tabellen aller Art erreichen lässt. Im selben Abschnitt findet sich ferner eine neue Methode, um die Bewegung auf Ellipsen und Hyperbeln, die sich der Parabel nähern, auf die Bewegung auf der Parabel zurückzuführen. Der zweite Abschnitt enthält ein Verfahren, um die geozentrischen Orte der Himmelskörper mithilfe dreier Koordinaten zu bestimmen. Der Verfasser hatte darüber bekanntlich schon früher eine Abhandlung vorgelegt. Hier wird das Verfahren noch wesentlich verbessert. Im selben Abschnitt werden auch die Differenzialänderungen des geozentrischen Orts aus den Differenzialänderungen der einzelnen Elemente bestimmt. Der dritte Abschnitt behandelt besonders ausführlich die äußerst wichtige Aufgabe, aus zwei Orten eines Himmelskörpers auf seiner Bahn die Elemente zu bestimmen. Vor allem die zweite Lösung gab Gelegenheit, sehr viele neue Beziehungen zu entwickeln. Um die praktische Brauchbarkeit der dargestellten Methoden deutlicher zu machen und auch weniger Geübten das Studium zu erleichtern, sind die wichtigeren Ergebnisse durchweg mit Beispielen erläutert, die aus wirklichen Fällen gewählt wurden und größtenteils miteinander zusammenhängen. Die Bewegung auf einer Hyperbel, die ja auch bei Kometen tatsächlich vorkommen kann, wird schon wegen ihres analytischen Interesses beinahe ebenso ausführlich behandelt wie die Bewegung auf einer Ellipse oder Parabel.

Die im ersten Teil dargestellten Untersuchungen sind gleichsam der Stoff, aus dem im zweiten Teil die Lösung der großen Aufgabe aufgebaut wird, durch die Analyse der geozentrischen Erscheinungen die Bahn des Himmelskörpers selbst zu bestimmen. Diese Arbeit stellt eigentlich mehrere sehr unterschiedliche Anforderungen. Eine Arbeit ist nötig, um die noch ganz unbekannte Bahn eines Himmelskörpers zum ersten Mal näherungsweise zu bestimmen. Eine andere ist nötig, um die schon näherungsweise bekannten Elemente anhand einer längeren, vielleicht viele Jahre umfassenden Beobachtungsreihe genauer auszuarbeiten. Die erste Aufgabe verlangt mehr Geschick, die andere mehr Arbeit. Für die erste wird man nur so viele Beobachtungen verwenden wie nötig, im Allgemeinen also drei. Die Aufgabe, aus drei vollständigen Beobachtungen die Bahn eines Himmelskörpers zu bestimmen, ist gewissermaßen die wichtigste Aufgabe
\SourceLeaf{64}des ganzen Werks und bildet den Inhalt des ersten Abschnitts des zweiten Teils. Die Lösung selbst lässt sich hier nicht in einem Auszug wiedergeben. Als Beleg für ihre breite Anwendbarkeit führen wir nur an, dass sie in den drei erläuternden Beispielen mit gleichem Erfolg auf Beobachtungen angewandt wird, die 22 Tage, 71 Tage beziehungsweise 260 Tage auseinanderliegen. Eine Besonderheit ist, dass man die kleinen Änderungen durch Parallaxe und Aberration nach Belieben \textit{sofort} berücksichtigen kann, ohne die Arbeit wesentlich zu erschweren. Ihretwegen braucht man die Rechnung also nicht doppelt auszuführen.

Unter gewissen Umständen, besonders wenn die Bahnebene nur wenig gegen die Ekliptik geneigt ist, wäre es nicht zweckmäßig, die Bestimmung der Elemente auf drei vollständige Beobachtungen zu stützen. Man muss vier Beobachtungen zugrunde legen, von denen aber nur zwei vollständig sein müssen. So ergibt sich eine neue Aufgabe, der der zweite Abschnitt gewidmet ist. Ein Beispiel zu Vesta, bei dem die erste und die letzte Beobachtung 162 Tage auseinanderliegen, erläutert diese Methode. Auch hier kann man, wenn man will, Parallaxe und Aberration \textit{sofort} berücksichtigen.

Zur genaueren Ausarbeitung der Elemente eines Himmelskörpers geht man gerade umgekehrt vor. Man verwendet alle verfügbaren Beobachtungen anstelle der kleinstmöglichen Zahl. Wie man dabei vorgeht, zeigt der dritte Abschnitt. Hier war der Ort, die Hauptpunkte einer für jede Anwendung der Mathematik auf die physische Welt höchst wichtigen Frage zu entwickeln. Wie verbindet man Beobachtungen und Messungen am zweckmäßigsten, um Ergebnisse zu gewinnen, wenn sie wegen der Unvollkommenheit unserer Sinne und Instrumente unvermeidlich mit Fehlern behaftet sind, mögen diese auch noch so klein sein? Die hier ausgeführten Grundsätze, die der Verfasser schon seit 14 Jahren angewandt und vor geraumer Zeit mehreren seiner astronomischen Freunde mitgeteilt hatte, führen zu der Methode, die auch \textsc{Legendre} vor einigen Jahren in seinem Werk \textit{Nouvelles méthodes pour la détermination des orbites des comètes} unter dem Namen \textit{Méthode des moindres carrés} aufgestellt hat. Die Begründung der Methode, die der Verfasser gibt, ist
\SourceLeaf{65}\SourceGathering{8\textsuperscript{*}}ganz seine eigene. Eine ausführlichere Behandlung ist von ihm in Zukunft zu erwarten.

Um eine \textit{bleibende} Übereinstimmung der Elemente mit den Bewegungen eines Planeten zu erreichen, darf man sich nicht auf eine rein elliptische Bewegung beschränken. Man muss auch die von den anderen Planeten verursachten Störungen berücksichtigen. Die Berechnung der Störungen selbst lag natürlich außerhalb des Plans eines Werks, das allein der Bewegung nach den Gesetzen \textsc{Kepler}s gewidmet ist. Der vierte und letzte Abschnitt deutet aber wenigstens die wichtigsten Schritte an, die nötig sind, um eine Planetenbahn unter Berücksichtigung der schon berechneten Störungen zu verbessern und ihr gleichsam den letzten Schliff zu geben.

Wir schließen diese Anzeige mit dem Wunsch, dass dieses Werk, die Frucht einer siebenjährigen Beschäftigung mit diesen Untersuchungen, dazu beitragen möge, diesem wichtigsten und schönsten Teil der theoretischen Astronomie mehr Freunde zu gewinnen und den Astronomen selbst ihre mit dem stetigen Wachstum der Wissenschaft immer umfangreicheren Arbeiten zu erleichtern.
\par\bigskip\centerline{\rule{.24\textwidth}{.3pt}}
